\documentclass[10pt]{article}
\usepackage[margin=20mm]{geometry}
\usepackage{amsmath,amssymb,fontspec,unicode-math,xurl,hyperref}
\IfFontExistsTF{DejaVu Serif}{\setmainfont{DejaVu Serif}}{\setmainfont{Latin Modern Roman}}
\setmathfont{latinmodern-math.otf}
\setlength{\parindent}{0pt}\setlength{\parskip}{6pt}\sloppy
\title{Scientific verification report
I25-40129 / 0027}
\author{}
\date{}
\begin{document}
\maketitle
Position: Medical Large Language Model Benchmarks Should Prioritize Construct Validity\par
Ahmed Alaa; Thomas Hartvigsen; Niloufar Golchini; Shiladitya Dutta; Frances Dean; Inioluwa Deborah Raji; Travis Zack\par
ICML 2025, PMLR 267:80991-81004; publisher-linked 14-page PDF\par
Minor: two bounded defects in the empirical argument; execution provenance unresolved\par
The documented EHR evaluation retains answer choices, so its performance drop does not isolate their removal. The text also names the wrong MedQA leader relative to Table 1. The recommendation to validate medical benchmarks and the observed ranking discordance are not invalidated.\par
\section*{1. Identity, scope and assumptions}
The preserved source [S1] is , SHA-256 `682e0ef017c005a8c69086ccaaba2e1cbe2eaab8d1d1fa6a26b61847df8bbd12`, matching the source-acquisition record. Title and authors were checked against the publisher record [S2]; the full middle name Deborah appears in that record and the PDF. Locations use 1-based PDF pages, also printed pages 1-14. Main text pp. 1-9 and Appendices A-C pp. 12-14 were inspected; references were used for discovery rather than exhaustively audited. Source pp. 6, 7, 8 and 12 were visually checked.\par
This is a position paper with a proof-of-concept EHR comparison. Its statistical targets include item-level conditional transfer accuracy, aggregate clinical-case accuracy, and agreement of rankings across models. These targets are distinct. Table 1 values are rounded published summaries, not patient-level data. No clinical records were accessed and no author code was run.\par
\section*{2. Problem 1: Answer-choice explanation is not isolated by the documented comparison}
Source claim and location. Section 4.1, p. 7, paragraph below Figure 5, proposes the contrived MedQA format as a likely explanation of the accuracy drop, specifically referring to answer choices being absent in real-world clinical reasoning. Appendix B, p. 12, instead describes pairing the retrieved EHR note, after removing its Plan section, with the extracted MedQA question and its answer choices [S1].\par
Complete check. Under the appendix protocol, the benchmark condition uses an MCQ clinical vignette, while the EHR condition uses an MCQ with a retrieved clinical context. Both retain answer choices. Therefore, the comparison changes context and case construction, but does not implement removal of choices. Its accuracy difference cannot identify the effect of a format intervention that was not performed. Figure 5’s separate answer-count experiment does not establish that the EHR comparison lacked choices.\par
Alternative readings and counterevidence. The main-text sentence can be read as a general ecological criticism of medical MCQs; that criticism is not refuted. It becomes unsupported when used to explain this particular accuracy drop. The appendix could describe a protocol different from the one actually used, but neither execution logs nor prompt records were available to decide that. The displayed measurements themselves have not been disproven.\par
Judgment and impact. Minor, confined to the methodological interpretation of the accuracy gap. The actual executed protocol is a Needs Evidence subitem, not an additional demonstrated error. If Appendix B is accurate, a supported repair is: ’The comparison shows a performance drop when MedQA questions and choices are paired with EHR-derived contexts; this design does not isolate the effect of removing answer choices.’ This is a conditional interpretation repair, not a validated experimental repair. A paired evaluation with and without choices would be new work.\par
\section*{3. Problem 2: The named MedQA leader conflicts with Table 1}
Source claim and location. Section 4.3, first paragraph on p. 8, identifies GPT-4 as the model at the top of the MedQA rankings. Table 1 on p. 6 lists the following relevant values [S1]:\par
Model | MedQA accuracy | EHR accuracy\par
PMY MedLLama | 0.75 | 0.36\par
GPT-4 | 0.71 | 0.28\par
Llama 3 | 0.54 | 0.48\par
Complete check. The maximum MedQA value among all seven displayed rows is 0.75 for PMY MedLLama, exceeding GPT-4 by 0.04. Llama 3 has the highest displayed EHR accuracy, 0.48 [E1]. These are direct finite comparisons, not significance tests. A pairwise comparison restricted to GPT-4 and Llama 3 would make the ordering correct, but that restriction is absent from the sentence referring to the full table.\par
Judgment, impact and repair. Minor, a local factual error in the model comparison used to motivate construct-validity concerns. The supported replacement is: ’PMY MedLLama leads Table 1 on MedQA (0.75), whereas Llama 3 leads the EHR-adapted evaluation (0.48).’ The numerical repair is verified against the table. Ranking discordance remains because PMY’s EHR accuracy is 0.36, below Llama 3’s 0.48. This repair does not establish that discordance uniquely diagnoses inadequate medical knowledge.\par
\section*{4. Checks that did not establish additional defects}
Abstention accounting. Table 1 reports that GPT-4 refrains on 57\% of cases and is correct on 64\% of answered cases. Counting abstentions as errors gives (1-0.57)*0.64 = 0.2752, consistent with 0.28 after rounding [E1]. Section 4.3 explicitly notes that nonresponse complicates interpretation. It would be incorrect to report this as an undisclosed arithmetic error.\par
Conditional transfer metric. Section 4.1 defines alpha = P(EHR correct | matched MedQA correct). This is a descriptive item-level measure; it is not interchangeable with across-model score prediction or rank agreement. For example, if model-specific benchmark success is theta and clinical success is theta/2, independently conditional on model, expected clinical score is exactly predicted by benchmark score, yet alpha equals theta/2 rather than one. This analytical distinction does not validate MedQA or refute the paper’s observed values. The alpha-near-one expectation on p. 7 is defensible only for the stronger target of near-deterministic item-wise transfer.\par
Clinical matching and content. The p. 7 discussion itself acknowledges that matching on a shared answer code may fail to capture clinical similarity. The actual note-question compatibility, adjudicated answers and repeated-patient structure remain unverified. Figure 5 and Appendix C were checked for described task categories, concept counts and perturbation scope, but the classifier outputs and cTAKES extractions were not reproduced. Document-level concept counts alone do not identify every aspect of construct coverage. These are scoped evidence limits, not demonstrated false empirical results.\par
\section*{5. Dependence, versions and decision table}
The two bounded findings affect the interpretation and named example of the proof of concept. They do not invalidate the normative proposal for benchmark validation, establish that MedQA is valid, or prove that the reported accuracies are wrong. Patient-level outputs, denominators, matching adjudications, model versions and complete prompts are needed for further statistical or clinical conclusions. Repeated notes per benchmark item also need to be retained when estimating uncertainty.\par
arXiv:2503.10694 lists v1 dated 12 March 2025 [S3]. That preprint was not substituted for the later publisher PDF. No erratum or linked author repository was identified in the inspected publisher/arXiv pages and targeted searches. OpenReview returned a browser-verification page [S4]; inaccessible discussion and attachments were not searched. Absence of a located correction is not proof that none exists.\par
Judgment / scope | Repair verification | Adoption / human status | Remaining work\par
Choice-format explanation: Minor | Conditional wording correction supported by Appendix B | Not requested / pending | Actual execution protocol needs evidence\par
Named MedQA leader: Minor | Published Table 1 supports PMY MedLLama | Not requested / pending | No re-estimation of raw accuracy\par
Refusal arithmetic: No Issue Found | 0.2752 correctly rounds to 0.28 | Not requested / pending | Published rounded inputs only\par
Metric scope / matching: clarification | No further defect or empirical repair established | Not requested / pending | Raw-data and construct-validity reanalysis\par
\section*{Supporting checks}
- Alaa et al. Position: Medical Large Language Model Benchmarks Should Prioritize Construct Validity. PMLR 267:80991-81004 (2025). Exact PDF: https://raw.githubusercontent.com/mlresearch/v267/main/assets/alaa25a/alaa25a.pdf . Preserved hash stated above.\par
- https://proceedings.mlr.press/v267/alaa25a.html . Snapshot: .\par
- arXiv:2503.10694v1, 12 March 2025. https://arxiv.org/abs/2503.10694 . Metadata snapshot in .\par
- OpenReview forum YuMEUNNpeb: https://openreview.net/forum?id=YuMEUNNpeb . Access returned browser verification; response archived as .\par
- and ; published rounded Table 1 values, ordering checks and refusal arithmetic.\par
- , and ; records retrievals and hashes.\par
Replay from verification/0027: python -X utf8 ; python -X utf8 ; python -X utf8 . Outputs: MedQA leader PMY MedLLama; EHR leader Llama 3; refusal-derived accuracy 0.2752; alpha-minus-marginal values range from 0.01 to 0.14. These are finite checks, not empirical reproduction.\par
Human review and scientific decisions remain pending. No author code or clinical data access was performed.\par
Scientific decisions and human verification remain pending; author adoption was not requested.\par
\end{document}
